\chapter{مقدمه}
ایران در مقطعی تاریخی قرار دارد که تصمیمات امروز، سرنوشت نسل‌های آینده را رقم خواهد زد. نه تعلل و بی‌عملی و یا ترس از تغییر کارساز است، نه شتاب و هیجان و احساساتی شدن و یا تمامیت‌خواهی و مصادره‌ی جنبش به نفع فرد یا گروهی خاص.

برای داشتن ایرانی آزاد، آباد، متکثر و چندصدایی، به شکلی که پایه‌های مستحکم انسجام و آینده‌نگری و توسعه‌ی پایدار در آن لحاظ شده باشد، باید فرایند گذار را هر چه بیشتر و بهتر و بر اساس دروس آموخته‌ی ملی و بین‌المللی مهندسی کرد.

خشم و خشونت و عمل‌گرایی سریع و انقلابی تا مقطع عبور از جمهوری اسلامی و اطمینان از بازگشت‌ناپذیر بودن آن قابل درک و شاید ضروری باشد — هرچند که هنوز نباید مصداق خشونت کور و یا خارج از اندازه باشد — ولی فراتر از آن، قطعاً و طبیعتاً مانع ساختن و همکاری و همدلی مردمی خواهد بود که همه زجر کشیده و آسیب‌دیده از این نظام قهرآمیز و سرکوب‌گر بوده و هستند. مردمی که از هر قشر و گروه و طبقه هزینه‌های مختلف و گزافی را پرداخت کرده و پرداخت می‌کنند و نمی‌توان این هزینه‌ها را محدود به برخی گروه‌ها و یا افکار و جریان‌ها کرد.

\begin{modernbox}{طراحی نهادی دوران گذار}
تجربه‌ی جهانی — از تونس و آفریقای جنوبی تا اسپانیا و لهستان و شیلی، و همچنین ضدالگوهایی مانند لیبی، مصر و عراق — نشان داده است که \textbf{کیفیت طراحی نهادی دوران گذار} بیش از هر عامل دیگری تعیین‌کننده‌ی سرنوشت نظام سیاسی آینده است. کشورهایی که با آگاهی، فراگیری و اجماع این مرحله را طی کرده‌اند، از دام‌های هرج‌ومرج، بازگشت استبداد و جنگ داخلی گریخته‌اند.
\end{modernbox}

\begin{center}
    \href{https://mahdisalem.com/documents/statements/Mahestan/Mahestan_2Pages_Summary.pdf}{
        \tcbox[colback=LightBlue, colframe=MainBlue, arc=5pt]{📥 دانلود خلاصه طرح مهستان (نسخه‌ی ۲ صفحه)}
    }
\end{center}
مهدی سالم؛ 1 مارچ 2026